\cotefile{123}
\scanpage{20}
nier (or, what amounts to the same thing, $c$) is well determined. The unipotents which define the said filtration, in terms of an $i$ which defines it, are \inserted{those} of the form $\exp(e)$, where $e$ ranges over $\handwritten{\mathfrak{g}}(2)^{\handwritten{*}}+\operatorname{Fil}^{\handwritten{3}}(\handwritten{\mathfrak{g}})$. They form a conjugacy class of elements of $U$ \uline{under} $P$ (NB but not under $U$), the conjugacy classes \uline{under} $U$ being identified with the sets of the form $(\exp(e)\handwritten{\mid}e\handwritten{\in}a+\operatorname{Fil}^{\handwritten{3}}(\handwritten{\mathfrak{g}},))$, where $a\handwritten{\in\mathfrak{g}}(2)^{\handwritten{*}}$ is an arbitrary element.\deleted{of}

\uline{Principal case}. For one to have\hspace{1.5em}$(1)=0$, it is necessary and sufficient that \deleted{the homomor} the class of homomorphisms $f:\operatorname{Sl}(2,k)\hspace{1.5em}G$ under consideration be ``\unsure{princi\deleted{a}p\deleted{p}al|princi\deleted{p}a\deleted{a}pl}'' in the sense of Dynkin-Kostant. Then the associated parabolic $P$ is a Borel (the converse not being true \deleted{probably}, cf Dynkin's table), and the group $Z_o$ is the trivial group (i.e. $\operatorname{Dyn}_c$ is a torsor under $G$), \deleted{\unsure{\$|S}}the converse probably not being true.
\margin{\handwritten{\uline{example}\quad Case $A_2,A_4$ the parabolics associated with the ``subregular'' unipotents are \unsure{Borels|borels} !}}
% archive sheet number: 19
